\documentclass[12pt]{article}
\usepackage[utf8]{inputenc}
\usepackage[T1]{fontenc}
\usepackage{hyperref}
\title{Thracian Language Resources for K-12 Curriculum Development}
\author{Compiled by Claude Sonnet 5 (claude-sonnet-5) for Andr\'as Kornai}
\date{}
\begin{document}
\maketitle

\section{Introduction}
Thracian (ISO 639-3 txh) is an extinct Indo-European language of the ancient Balkans. It has no native speakers and no continuous community of heritage speakers. No text has been satisfactorily deciphered, so the language is known only from a small number of inscriptions, Greek glosses and personal and place names. Nothing in my searches suggests that a school curriculum in Thracian exists or could be built on fluent speakers. Any K-12 use would have to be about the history and heritage of the language, not instruction in it.

\section{Primary Educational Resources}
I found no primers, textbooks or school materials for Thracian.

\section{Linguistic and Reference Materials}
\begin{itemize}
\item Overview article on the Thracian language (Wikipedia): \url{https://en.wikipedia.org/wiki/Thracian_language}
\item Corpus of Inscriptions in the Local Language(s) of Thrace. This is a project at Kiel University that plans an open-access digital edition in EpiDoc format. A project report gives about 400 inscriptions plus glosses. Abstracts: \url{https://epigraphy.info/documents/workshop_10/abstracts/Santamaria-Selvi.html} and \url{https://epigraphy.info/documents/workshop_10/abstracts/Santamaria.html}. Report: \url{https://www.ku.de/fileadmin/140201/Ernst-Kirsten-Gesellschaft/Report_EKG_Santamaria.pdf}
\item Dan Dana, chapter on inscriptions from the Thracian area, which includes a reconsideration of texts written in Thracian (free PDF on HAL): \url{https://halshs.archives-ouvertes.fr/halshs-01405153}
\item Cambridge Scholars Publishing sample chapter: \url{https://www.cambridgescholars.com/resources/pdfs/978-1-4438-1325-9-sample.pdf}
\item Mircheva-Zhechkova on the database of Thracian glosses: \url{https://uni-sofia.academia.edu/AlbenaMirchevaZhechkova}
\end{itemize}
Classic printed works that the search results mention but do not link are Detschew, \emph{Die thrakischen Sprachreste}, and Dana, \emph{Onomasticon Thracicum} (2014).

\section{Religious and Community Resources}
I found no religious or community-language resources.

\section{Digital and Audio Resources}
I found no audio or software resources. The Kiel corpus is planned as an open-access digital tool, but I did not confirm that it is online yet.

\section{Curriculum Development Context}
Thracian is attested only fragmentarily. A language-medium curriculum is not possible. A heritage unit could use the inscriptions, glosses and names as classroom material on Balkan history and historical linguistics. The Kiel corpus is the most promising source once it is published.
\end{document}